\section{统计与线性代数}

\answer{AG1-1}
由 $\Cov(W,Z)=\Var(T)$ 得潜在方差为 $3$。因此 $\Var(U)=5-3=2$、$\Var(V)=4-3=1$，均非负。工具变量比值给出 $\beta=6/3=2$，而普通回归斜率为 $6/5$。取三个相互独立、方差分别为 $3,2,1$ 的中心正态变量，并取独立中心响应噪声，即可实现这些给定矩。
该模型必有 $\Cov(W,Y)=\Cov(Z,Y)=\beta\Var(T)$，所以若两协方差变为 $6$ 和 $5$，则彼此矛盾；不能通过改变测量噪声方差来消除矛盾。这是模型的可检验限制，不能任意挑一个比值当作正确答案。

\answer{AG1-2}
A25 给出两者共同的界 $1/(1+2^2)=1/5$。非严格事件在 $X=2$ 的概率为 $1/5$、$X=-1/2$ 的概率为 $4/5$ 时达到。均值为 $2/5-2/5=0$，二阶矩为 $4/5+1/5=1$。
严格事件不能达到该界，否则同时达到非严格界，迫使分布恰为上述两点分布，而其严格尾概率为零。对任意 $h>0$，令 $b=2+h$，取 $X=b$ 的概率为 $1/(1+b^2)$，取 $X=-1/b$ 的概率为 $b^2/(1+b^2)$。均值零、方差一，严格尾概率随 $h\downarrow0$ 趋于 $1/5$。

\answer{AG1-3}
省略常数后目标为 $\sum_j(b_j-z_j)^2/2+\sum_jb_j/2$。给预算约束乘子 $\nu\ge0$，固定它时非负约束下各坐标的最优值为
\[
 b_j=(z_j-1/2-\nu)_+.
\]
若 $\nu=0$，得到 $(5/2,1/2,0)$，总和为 $3$，不满足预算。因此预算必须取等。取 $\nu=1/2$ 得唯一解 $(2,0,0)$，总和恰为 $2$。第一坐标梯度加乘子为零，第二坐标在零点的单侧导数为零，第三为正，故满足全部最优性条件。二次项严格凸，所以即使第二坐标处于阈值边界，解仍唯一。预算约束通过共同乘子耦合了各个坐标。

\answer{AG1-4}
取 $c=(1,1)^\top,d=0$。A06 中 $e=1$、$q=1+1/4=5/4$，所以
\[
 \widehat\beta_c=(2,-1)^\top-\frac45(1,1/4)^\top
 =(6/5,-6/5)^\top,
 \qquad\Delta\mathrm{RSS}=4/5.
\]
直接令 $\beta=(t,-t)$，额外平方和为 $(t-2)^2+4(1-t)^2$。其导数为 $10t-12$，在 $t=6/5$ 唯一为零；代入得 $16/25+4/25=4/5$。这也说明对两个系数做等量校正并非本题的最优做法。

\par\Needspace{8\baselineskip}
\subsection*{小节练习}
\answer{OA02}
中心化不改变斜率，以下设各变量均值为零。记联合回归为
\[
 Y=bX_1+\beta_2X_2+u,
 \qquad \Cov(u,X_1)=\Cov(u,X_2)=0.
\]
令 $\gamma=\Cov(X_1,X_2)/\Var(X_1)$，并将 $X_2$ 对 $X_1$ 的残差记为 $r=X_2-\gamma X_1$。由联合回归的第一条正规方程，$Y$ 对 $X_1$ 的单变量斜率为 $b+\beta_2\gamma$，故
\[
 e=\beta_2r+u.
\]
由于 $r$ 与 $X_1$ 不相关，且 $X_2=r+\gamma X_1$，有
\[
 \Cov(r,X_2)=\Var(r)
 =\Var(X_2)(1-\rho^2).
\]
因此
\[
 \boxed{\beta_1
 =\frac{\Cov(e,X_2)}{\Var(X_2)}
 =(1-\rho^2)\beta_2.}
\]
只对响应做残差化时，第二次回归的分母仍为 $\Var(X_2)$。若同时对解释变量做残差化，即回归 $e$ 对 $r$，分母变为 $\Var(r)$，所得斜率才是 $\beta_2$。当 $\rho=0$ 时两系数相同；$|\rho|=1$ 时联合回归的两个斜率一般不能分别识别，因此题目排除了这一情形。

\answer{OA03}
记 $A=\{X+Y>1\}$。原来的联合密度在单位正方形上恒为 $1$，而 $\Prob(A)=1/2$，故条件密度在上三角形内恒为 $2$。给定 $X=x$ 与 $A$ 后，$Y$ 在 $(1-x,1)$ 上均匀分布，于是
\[
 \E[Y\mid X=x,A]=1-\frac x2,
 \qquad 0<x<1.
\]
条件均值已经是仿射函数。令 $u=Y-(1-X/2)$，则 $\E[u\mid X,A]=0$，从而 $u$ 与常数及 $X$ 正交。因此最小二乘回归的截距为 $1$、斜率为
\[
 \boxed{-\frac12}.
\]
也可用矩核对。条件密度 $f_{X\mid A}(x)=2x$ 给出
\[
 \E[X\mid A]=\frac23,\quad
 \E[X^2\mid A]=\frac12,\quad
 \Var(X\mid A)=\frac1{18}.
\]
由条件均值公式，$\Cov(X,Y\mid A)=-\Var(X\mid A)/2=-1/36$，协方差除以方差仍得到 $-1/2$。独立性属于筛选前的总体，不能直接用于条件总体。

\answer{OA05}
令 $c=X^\top y/n$。展开目标函数之差得
\[
 L(\beta)-L(0)
 =\frac{\|X\beta\|^2}{2n}-c^\top\beta
   +\lambda\|\beta\|_1.
\]
若 $\|c\|_\infty\le\lambda$，由 $c^\top\beta\le\|c\|_\infty\|\beta\|_1$，右侧非负，故零向量最优。
反之，若某个坐标满足 $|c_j|>\lambda$，令 $\beta=t\operatorname{sign}(c_j)e_j$，其中 $t>0$。此时
\[
 L(\beta)-L(0)
 =\frac{t^2\|X_j\|^2}{2n}-t(|c_j|-\lambda).
\]
对充分小的 $t$，负的一次项支配二次项，目标值严格下降。因此充要条件为
\[
 \boxed{\lambda\ge\frac{\|X^\top y\|_\infty}{n}}.
\]
在该条件下，若 $\lambda>0$，零解还必然唯一：若另一最优解满足 $X\beta\ne0$，损失对拟合值的严格凸性会使它与零解的中点取得更小目标值；若 $X\beta=0$ 而 $\beta\ne0$，惩罚项严格增加。若 $\lambda=0$，条件化为 $X^\top y=0$，此时所有最优解恰为 $\ker X$，只有当 $X$ 满列秩时零解才唯一。

\answer{OA30}
结论为
\[
 \boxed{\lambda_i\le\mu_i\le\lambda_{i+1}},
 \qquad i=1,\ldots,n-1.
\]
令 $H=\{x\in\R^n:x_k=0\}$。把 $B$ 的向量在第 $k$ 个坐标补零后，$B$ 的 Rayleigh 商就是 $A$ 的 Rayleigh 商在 $H$ 上的限制。记 $R_A(x)=x^\top Ax/\|x\|^2$，由极小极大原理，
\[
 \lambda_i=\min_{\dim S=i}\ \max_{0\ne x\in S}R_A(x),
\]
而 $\mu_i$ 的同一极小值只允许取 $S\subset H$。可选子空间变少，故 $\mu_i\ge\lambda_i$。

再令 $E$ 为 $A$ 前 $i+1$ 个特征向量张成的空间。由于 $H$ 的余维为一，$\dim(E\cap H)\ge i$。从交空间取一个 $i$ 维子空间 $S$，其上所有 Rayleigh 商均不超过 $\lambda_{i+1}$。把它代入 $\mu_i$ 的极小极大公式，即得 $\mu_i\le\lambda_{i+1}$。

特征值可以重复，因此不等号不能统一改为严格不等号。例如 $A$ 为单位矩阵时，两侧均取等。

